%! Author = kristofferpaulsson
%! Date = 2024-08-03

% Preamble
\documentclass[10pt,a4paper,notitlepage]{article}

\usepackage[margin=2.5cm]{geometry}
\usepackage{fontspec}
\usepackage[none]{hyphenat}  % --- NO HYPHENATION

\usepackage[hyphens]{xurl}  % --- URL HYPHENATION
\usepackage{hyperref}  % --- URL

\usepackage{draftwatermark}  % --- DONT-CITE

\urlstyle{same}  % --- URL
\pagenumbering{gobble}
\DraftwatermarkOptions{fontsize=2em,angle=0,vpos=1cm,color=black,text=Do Not Cite or Circulate}  % --- DONT-CITE

\usepackage{float}  % --- TABLE
\usepackage{makecell}  % --- TABLE
\usepackage{enumitem}  % --- LIST


\fontspec{Tinos}[
    Path = ../../font/t/,
    Extension = .ttf,
    UprightFont = *-Regular,
    Scale=1,
    Ligatures=TeX
]
\fontspec{Inter}[
    Path = ./../font/i/,
    Extension = .ttf ,
    UprightFont = *-Regular ,
    Scale=1,
    Ligatures=TeX
]

% Commands
\newcommand{\tsc}[1]{\textsc{#1}}
\newcommand*{\grc}[1]{\fontspec{Inter}#1\rmfamily}
\newcommand{\trc}[1]{\textit{\fontspec{Tinos}#1}}
\newcommand{\linkfoot}[3]{\footnote{\emph{#1}, s.v. ``{#2},'' accessed \printdate{#3}.}}

% Document
\begin{document}

    \subsection*{Ancient Greek: Indefinite Relative Pronouns}
    \begin{itemize}[label={}, leftmargin=0]
        \item Kristoffer E.D.\ Paulsson (2026), Independent Scholar.
    \end{itemize}

    ``The indefinite or general relative pronoun \grc{ὅστις}\fontspec{}, \grc{ἥτις}\fontspec{}, \grc{ὅτι} \fontspec{} whoever (any-who, any-which), \emph{any one who, whatever}, \textit{anything which}, inflects each part (\grc{ὅς} \fontspec{} and \grc{τὶς}\fontspec{}) separately.''

    ``and \grc{ὅστις}, \grc{ἥτις}, \grc{ὅ τι}, \textit{whoever, whichever.} The latter is called the indefinite relative.''

    %The following table contains indefinite relative pronouns described with natural gender, a common use is grammatical gender, which currently isn't included.
    %Therefore the word man/men refers to male agent(s), and woman/women refers to a female agent(s).

    %Greek pronouns used in Homeric and other kind of prose excluded from the table.

    \begin{table}[H]
        \begin{tabular}{c|ccc|c}
            & \multicolumn{3}{c|}{Sexus} & Genus \\
            & \tsc{Masc} & \tsc{Fem} & \tsc{Neu} & \tsc{M. F. N.} \\
            \hline
            \emph{\tsc{Singular}} & & & & \\
            \tsc{Nom} & \makecell{\grc{ὅστις} \trc{hostis} \\ \small any man who} & \makecell{\grc{ἥτις} \trc{hētis} \\ \small any woman who } & \makecell{\grc{ὅτι} \trc{hoti} \\ \small whoever} & \makecell{( * ) \\ \small whatever} \\
            \tsc{Gen} & \makecell{\grc{οὗτινος} \trc{houtinos} \\ \small any man whose} & \makecell{\grc{ἧστινος} \trc{hēstinos} \\ \small any woman whose} & \makecell{\grc{οὗτινος} \trc{houtinos} \\ \small whosever } & \makecell{( * ) \\ \small of whatever} \\
            \tsc{Dat} & \makecell{\grc{ᾧτινι} \trc{hōitini} \\ \small any man to whom } & \makecell{\grc{ᾗτινι} \trc{hēitini} \\ \small any woman to whom} & \makecell{\grc{ᾧτινι} \trc{hōitini} \\ \small to whomever} & \makecell{( * ) \\ \small to whatever} \\
            \tsc{Acc} & \makecell{\grc{ὅντινα} \trc{hontina} \\ \small any man whom} & \makecell{\grc{ἥντινα} \trc{hēntina} \\ \small any woman whom} & \makecell{\grc{ὅτι} \trc{hoti} \\ \small whomever} & \makecell{( * ) \\ \small whatever} \\
            \hline
            \emph{\tsc{Dual}} & & & & \\
            \tsc{Nom/Acc} & \makecell{\grc{ὥτινε} \trc{hōtine} \\ \small any two men who(m)} & \makecell{\grc{ὥτινε} \trc{hōtine} \\ \small any two women who(m)} & \makecell{\grc{ὥτινε} \trc{hōtine} \\ \small any two who(m)} & \makecell{( * ) \\ \small that which} \\
            \tsc{Dat/Gen} & \makecell{\grc{οἷντινοιν} \trc{hointinoin} \\ \small any two men \\ \small whose/to whom} & \makecell{\grc{οἷντινοιν} \trc{hointinoin} \\ \small any two women \\ \small whose/to whom} & \makecell{\grc{οἷντινοιν} \trc{hointinoin} \\ \small any two who(m/se)} & \makecell{( * ) \\ \small that which} \\
            \hline
            \emph{\tsc{Plural}} & & & & \\
            \tsc{Nom} & \makecell{\grc{οἵτινες} \trc{hoitines} \\ \small any men who} & \makecell{\grc{αἵτινες} \trc{haitines} \\ \small any women who} & \makecell{\grc{ἅτινα} \trc{hatina} \\ \small those who} & \makecell{( * ) \\ \small that which} \\
            \tsc{Gen} & \makecell{\grc{ὧντινων} \trc{hōntinōn} \\ \small any men whose} & \makecell{\grc{ὧντινων} \trc{hōntinōn} \\ \small any women whose} & \makecell{\grc{ὧντινων} \trc{hōntinōn} \\ \small those whose} & \makecell{( * ) \\ \small that which} \\
            \tsc{Dat} & \makecell{\grc{οἷστισι(ν)} \trc{hoistisi(n)} \\ \small any men to whom} & \makecell{\grc{αἷστισι(ν)} \trc{haistisi(n)}\\ \small any women to whom} & \makecell{\grc{οἷστισι(ν)} \trc{hoistisi(n)} \\ \small those whom} & \makecell{( * ) \\ \small that which} \\
            \tsc{Acc} & \makecell{\grc{οὕστινας} \trc{houstinas} \\ \small any men whom} & \makecell{\grc{ἅστινας} \trc{hastinas} \\ \small any women whom} & \makecell{\grc{ἅτινα} \trc{hatina} \\ \small those whom} & \makecell{( * ) \\ \small that which} \\
        \end{tabular}
        \caption{Compiled from standard Ancient Greek grammars enhanced with transliteration and cues.
            Assume gender over sex except when the antecedent is a person.
            (Smyth, 1956:97, \S339)}
    \end{table}

    \subsection*{Reference List}
    \begin{itemize}[label={}, leftmargin=0]
        \item Rydberg-Cox, J. (n.d.). \textit{LESSON LVI: Relative Pronouns. Genitive Absolute. Numerals}. [online] A Digital Tutorial for Ancient Greek. Available at: \url{https://daedalus.umkc.edu/FirstGreekBook/JWW_FGB56.html} [Accessed 6 Feb. 2026].
        %\item Mounce, W.D. (2013). \textit{Greek for the Rest of Us: the Essentials of Biblical Greek}. Grand Rapids, Michigan: Zondervan.
        \item Smyth, H.W. (1956). \textit{Greek Grammar}. Cambridge, Massachusetts: Harvard University Press.
    \end{itemize}

\end{document}
